% GENERATED FILE. Edit canonical lessons, then run npm run generate:books.
% canonical-source-hash: fnv1a64:c4d757b8d450cc9d
% canonical-lessons: KA-C209-manassu, KA-C209-nirase, KA-C209-tamase, KA-C209-yoga, KA-C209-dhyana

\chapter{Mind, Disappointment, Fun, Yoga, Meditation}
\label{ch:ka-mind-disappointment-fun-yoga-meditation}
\hlchaptermodality{209}

\begin{chapteropening}
\textbf{By the end of this chapter:} \emph{I can say the mind, disappointment, fun, yoga and meditation.}

Complete the last lesson of chapter 209: I can say the mind, disappointment, fun, yoga and meditation.
\end{chapteropening}

\section[\texorpdfstring{\kn{ಮನಸ್ಸು} (manassu)}{ಮನಸ್ಸು (manassu)}]{\kn{ಮನಸ್ಸು} (manassu) — the mind}

\label{lesson:KA-C209-manassu}



\begin{quote}
Before the new one: say the Kannada for an illness, then the Kannada for a patient.
\end{quote}

\subsection*{You'll want to know: \kn{ಮನಸ್ಸು}}
\textbf{\kn{ಮನಸ್ಸು}} — \emph{manassu} — \textquotedblleft{}the mind\textquotedblright{}.

\begin{grammarlens}[title={What you've built}]
One more word for this chapter.
\end{grammarlens}

\subsection*{Your turn}
\begin{itemize}
\raggedright
  \item \emph{Say it:} \emph{manassu}
  \item \emph{Say it:} \emph{manassu}, once more
  \item \emph{Say it:} say \emph{manassu}
  \item \emph{Recall:} say the Kannada for an illness, then the Kannada for a patient, then say \emph{manassu} again
\end{itemize}

\begin{tcolorbox}[breakable,colback=teal!4,colframe=teal!35!black,title={Before you move on}]
What does \kn{ಮನಸ್ಸು} mean? (\textquotedblleft{}The mind\textquotedblright{}.) Say it once more.
\end{tcolorbox}
\section[\texorpdfstring{\kn{ನಿರಾಸೆ} (nirāse)}{ನಿರಾಸೆ (nirāse)}]{\kn{ನಿರಾಸೆ} (nirāse) — disappointment}

\label{lesson:KA-C209-nirase}



\begin{quote}
Before the new one: say the Kannada for a patient, then the Kannada for the mind.
\end{quote}

\subsection*{You'll want to know: \kn{ನಿರಾಸೆ}}
\textbf{\kn{ನಿರಾಸೆ}} — \emph{nirāse} — \textquotedblleft{}disappointment\textquotedblright{}.

\begin{grammarlens}[title={What you've built}]
One more word for this chapter.
\end{grammarlens}

\subsection*{Your turn}
\begin{itemize}
\raggedright
  \item \emph{Say it:} \emph{nirāse}
  \item \emph{Say it:} \emph{nirāse}, once more
  \item \emph{Say it:} say \emph{nirāse}
  \item \emph{Recall:} say the Kannada for a patient, then the Kannada for the mind, then say \emph{nirāse} again
\end{itemize}

\begin{tcolorbox}[breakable,colback=teal!4,colframe=teal!35!black,title={Before you move on}]
What does \kn{ನಿರಾಸೆ} mean? (\textquotedblleft{}Disappointment\textquotedblright{}.) Say it once more.
\end{tcolorbox}
\section[\texorpdfstring{\kn{ತಮಾಷೆ} (tamāṣe)}{ತಮಾಷೆ (tamāṣe)}]{\kn{ತಮಾಷೆ} (tamāṣe) — fun, a joke}

\label{lesson:KA-C209-tamase}



\begin{quote}
Before the new one: say the Kannada for the mind, then the Kannada for disappointment.
\end{quote}

\subsection*{You'll want to know: \kn{ತಮಾಷೆ}}
\textbf{\kn{ತಮಾಷೆ}} — \emph{tamāṣe} — \textquotedblleft{}fun, a joke\textquotedblright{}.

\begin{grammarlens}[title={What you've built}]
One more word for this chapter.
\end{grammarlens}

\subsection*{Your turn}
\begin{itemize}
\raggedright
  \item \emph{Say it:} \emph{tamāṣe}
  \item \emph{Say it:} \emph{tamāṣe}, once more
  \item \emph{Say it:} say \emph{tamāṣe}
  \item \emph{Recall:} say the Kannada for the mind, then the Kannada for disappointment, then say \emph{tamāṣe} again
\end{itemize}

\begin{tcolorbox}[breakable,colback=teal!4,colframe=teal!35!black,title={Before you move on}]
What does \kn{ತಮಾಷೆ} mean? (\textquotedblleft{}Fun, a joke\textquotedblright{}.) Say it once more.
\end{tcolorbox}
\section[\texorpdfstring{\kn{ಯೋಗ} (yōga)}{ಯೋಗ (yōga)}]{\kn{ಯೋಗ} (yōga) — yoga}

\label{lesson:KA-C209-yoga}



\begin{quote}
Before the new one: say the Kannada for disappointment, then the Kannada for fun.
\end{quote}

\subsection*{You'll want to know: \kn{ಯೋಗ}}
\textbf{\kn{ಯೋಗ}} — \emph{yōga} — \textquotedblleft{}yoga\textquotedblright{}.

\begin{grammarlens}[title={What you've built}]
One more word for this chapter.
\end{grammarlens}

\subsection*{Your turn}
\begin{itemize}
\raggedright
  \item \emph{Say it:} \emph{yōga}
  \item \emph{Say it:} \emph{yōga}, once more
  \item \emph{Say it:} say \emph{yōga}
  \item \emph{Recall:} say the Kannada for disappointment, then the Kannada for fun, then say \emph{yōga} again
\end{itemize}

\begin{tcolorbox}[breakable,colback=teal!4,colframe=teal!35!black,title={Before you move on}]
What does \kn{ಯೋಗ} mean? (\textquotedblleft{}Yoga\textquotedblright{}.) Say it once more.
\end{tcolorbox}
\section[\texorpdfstring{\kn{ಧ್ಯಾನ} (dhyāna)}{ಧ್ಯಾನ (dhyāna)}]{\kn{ಧ್ಯಾನ} (dhyāna) — meditation}

\label{lesson:KA-C209-dhyana}



\begin{quote}
Before the new one: say the Kannada for fun, then the Kannada for yoga.
\end{quote}

\subsection*{You'll want to know: \kn{ಧ್ಯಾನ}}
\textbf{\kn{ಧ್ಯಾನ}} — \emph{dhyāna} — \textquotedblleft{}meditation\textquotedblright{}.

\begin{grammarlens}[title={What you've built}]
One more word for this chapter.
\end{grammarlens}

\subsection*{Your turn}
\begin{itemize}
\raggedright
  \item \emph{Say it:} \emph{dhyāna}
  \item \emph{Say it:} \emph{dhyāna}, once more
  \item \emph{Say it:} say \emph{dhyāna}
  \item \emph{Recall:} say the Kannada for fun, then the Kannada for yoga, then say \emph{dhyāna} again
\end{itemize}

\begin{tcolorbox}[breakable,colback=teal!4,colframe=teal!35!black,title={Before you move on}]
What does \kn{ಧ್ಯಾನ} mean? (\textquotedblleft{}Meditation\textquotedblright{}.) Say it once more.
\end{tcolorbox}
